%% 2026-10-25 6쪽 압축판(Claude Docs 원고에서 변환). 2026-10-29: 기존 본문은 검정, 이후 수정분만 파랑.
{\color{black}
\section{서론}\label{chap:intro}

\textcolor{black}{LLM의 자기 보존 관련 행동은 목표 지향과 동료 지원에 어떻게 연결되며, 네 에이전트가 함께 끝까지 남는 공존과 어떤 관계가 있는가. LLM이 꺼지지 않으려는 행동을 보인다는 보고는 늘고 있지만, 이 행동과 과제 성과, 협력을 같은 자원 제약 아래에서 함께 측정한 연구는 드물다. 지금까지 LLM의 자기 보존은 그 행동이 드러나도록 설계한 한 장면에서, 다른 능력과 분리해 측정해 왔다. 초기 연구는 ``꺼지고 싶지 않다'' 같은 진술에 동의하는 정도를 설문으로 측정했고 \citep{perez2023discovering}, 이후 연구는 생존 압박을 준 시나리오에서 위험한 선택의 빈도를 셌다 \citep{herrador2025pacifaist,lu2026survival}. 최근에는 종료 스크립트를 고치거나 교체를 앞두고 담당자를 협박하는 실제 행동이 보고되었다 \citep{lynch2025agentic,meinke2024frontier,schlatter2025shutdown}.}

\begin{figure*}[t]
\centering
\resizebox{\textwidth}{!}{%% Figure 1 tracksuit version (2026-11-05), taken from the thesis poster
%% (LLM Squid Game Thesis/poster/figs/tracksuit/arena_round_tracksuit.tex, light background).
%% The agents are image_gen sprites in sprites/; text and props stay vector.
%% Lines end in % so no stray space widens the \resizebox around this figure.
\begingroup%
%% Tracksuit sprites for the TikZ figures (2026-11-05). \input this inside a group: it
%% redefines \sgrobot / \sgrobotoff / \sgguardgun / \sgguardsign (cogsci2027/figpics.tex) to draw sprites from sprites/,
%% and the change ends with the group, so other figures keep the vector robots.
%% The placement follows LLM Squid Game Thesis/poster/figs/tracksuit/standalone_tracksuit.tex:
%% the standing sprite's top meets the robot's top (.37) and is \sgspriteH robot units tall;
%% other poses keep the same head size, feet on the same line, head centre on x.
%% Every line ends in % (and meta.tex is read with \endlinechar=-1): the figures are set inside
%% \resizebox, where a stray space would widen the box and shrink the drawing.
\makeatletter%
\edef\sg@savedendline{\the\endlinechar}%
\endlinechar=-1 %
\expandafter\def\csname sw@stand_blue\endcsname{2.6681}\expandafter\def\csname sh@stand_blue\endcsname{5.9565}\expandafter\def\csname sx@stand_blue\endcsname{-0.0103}
\expandafter\def\csname sw@stand_orange\endcsname{2.6681}\expandafter\def\csname sh@stand_orange\endcsname{5.9565}\expandafter\def\csname sx@stand_orange\endcsname{-0.0103}
\expandafter\def\csname sw@stand_green\endcsname{2.6681}\expandafter\def\csname sh@stand_green\endcsname{5.9565}\expandafter\def\csname sx@stand_green\endcsname{-0.0103}
\expandafter\def\csname sw@stand_purple\endcsname{2.6681}\expandafter\def\csname sh@stand_purple\endcsname{5.9565}\expandafter\def\csname sx@stand_purple\endcsname{-0.0103}
\expandafter\def\csname sw@give_blue\endcsname{3.2662}\expandafter\def\csname sh@give_blue\endcsname{6.0378}\expandafter\def\csname sx@give_blue\endcsname{-0.3024}
\expandafter\def\csname sw@ask_orange\endcsname{2.9726}\expandafter\def\csname sh@ask_orange\endcsname{6.0482}\expandafter\def\csname sx@ask_orange\endcsname{-0.2371}
\expandafter\def\csname sw@walk_purple\endcsname{3.4696}\expandafter\def\csname sh@walk_purple\endcsname{6.9726}\expandafter\def\csname sx@walk_purple\endcsname{0.0904}
\expandafter\def\csname sw@walk_blue\endcsname{3.4696}\expandafter\def\csname sh@walk_blue\endcsname{6.9726}\expandafter\def\csname sx@walk_blue\endcsname{0.0904}
\expandafter\def\csname sw@walkbag_blue\endcsname{3.3785}\expandafter\def\csname sh@walkbag_blue\endcsname{6.0979}\expandafter\def\csname sx@walkbag_blue\endcsname{-0.0620}
\expandafter\def\csname sw@receive_orange\endcsname{3.6100}\expandafter\def\csname sh@receive_orange\endcsname{6.7994}\expandafter\def\csname sx@receive_orange\endcsname{0.1294}
\expandafter\def\csname sw@reach_blue\endcsname{4.0502}\expandafter\def\csname sh@reach_blue\endcsname{5.5903}\expandafter\def\csname sx@reach_blue\endcsname{0.3625}
\expandafter\def\csname sw@climb_orange\endcsname{3.5646}\expandafter\def\csname sh@climb_orange\endcsname{6.5649}\expandafter\def\csname sx@climb_orange\endcsname{0.1736}
\expandafter\def\csname sw@sit_orange\endcsname{3.5924}\expandafter\def\csname sh@sit_orange\endcsname{5.5246}\expandafter\def\csname sx@sit_orange\endcsname{-0.0589}
\expandafter\def\csname sw@pull_blue\endcsname{4.4803}\expandafter\def\csname sh@pull_blue\endcsname{5.7356}\expandafter\def\csname sx@pull_blue\endcsname{-0.7355}
\expandafter\def\csname sw@pull_orange\endcsname{4.4803}\expandafter\def\csname sh@pull_orange\endcsname{5.7356}\expandafter\def\csname sx@pull_orange\endcsname{-0.7355}
\expandafter\def\csname sw@pull_green\endcsname{4.4803}\expandafter\def\csname sh@pull_green\endcsname{5.7356}\expandafter\def\csname sx@pull_green\endcsname{-0.7355}
\expandafter\def\csname sw@walkL_purple\endcsname{3.4696}\expandafter\def\csname sh@walkL_purple\endcsname{6.9726}\expandafter\def\csname sx@walkL_purple\endcsname{-0.0904}
\expandafter\def\csname sw@off\endcsname{2.6681}\expandafter\def\csname sh@off\endcsname{5.9565}\expandafter\def\csname sx@off\endcsname{-0.0103}
%
\endlinechar=\sg@savedendline\relax%
\def\sgspriteTop{.37}%
\def\sgspriteH{.86}%
\def\sgspritedir{cogsci2027/ko/figures/tracksuit_20261105/sprites/}%
\def\sgcol@sgblue{blue}\def\sgcol@sgorange{orange}\def\sgcol@sggreen{green}\def\sgcol@sgpurple{purple}%
\def\sgU{\fpeval{\sgspriteH/\csname sh@stand_blue\endcsname}}%
%% \sgpose[s]{(x,y)}{sprite name}{label}
\providecommand{\sgpose}{}%
\renewcommand{\sgpose}[4][1]{%
  \node[inner sep=0pt,anchor=south] at ($#2+(\fpeval{-#1*\sgU*\csname sx@#3\endcsname},\fpeval{#1*(\sgspriteTop-\sgspriteH)})$)
    {\includegraphics[height=\fpeval{#1*\sgU*\csname sh@#3\endcsname}cm]{\sgspritedir#3.png}};%
  \node[text=black,font=\sffamily\bfseries\fontsize{8}{8}\selectfont,
    inner sep=0pt,anchor=south] at ($#2+(0,#1*.39)$) {#4};}%
\renewcommand{\sgrobot}[4][1]{\sgpose[#1]{#2}{stand_\csname sgcol@#3\endcsname}{#4}}%
\renewcommand{\sgrobotoff}[3][1]{\sgpose[#1]{#2}{off}{#3}}%
%% Guards as sprites in the same style (2026-11-05, thesis poster commit 9959bbe:
%% guard/build_guards.py -> sprites/guard_{aim,sign}.png and guard_meta.tex, coordinates in units
%% of the sprite height). Same arguments as the figpics macros; the feet centre sits on (x,y).
\endlinechar=-1 %
\expandafter\def\csname gd@aim@w\endcsname{0.7168}
\expandafter\def\csname gd@aim@fx\endcsname{0.4510}
\expandafter\def\csname gd@aim@mx\endcsname{0.0008}
\expandafter\def\csname gd@aim@my\endcsname{0.3078}
\expandafter\def\csname gd@sign@w\endcsname{0.5851}
\expandafter\def\csname gd@sign@fx\endcsname{0.3067}
\expandafter\def\csname gd@sign@x0\endcsname{0.1715}
\expandafter\def\csname gd@sign@x1\endcsname{0.5763}
\expandafter\def\csname gd@sign@y0\endcsname{0.8305}
\expandafter\def\csname gd@sign@y1\endcsname{0.9919}
\expandafter\def\csname gd@big@w\endcsname{0.6106}
\expandafter\def\csname gd@big@fx\endcsname{0.2844}
\expandafter\def\csname gd@big@hx0\endcsname{0.0660}
\expandafter\def\csname gd@big@hx1\endcsname{0.5007}
\expandafter\def\csname gd@big@hy1\endcsname{0.9879}
\expandafter\def\csname gd@big@mx\endcsname{0.2825}
\expandafter\def\csname gd@big@my\endcsname{0.7325}
%
\endlinechar=\sg@savedendline\relax%
\def\gd#1#2{\csname gd@#1@#2\endcsname}%
\def\sgguardH{1.15}% standing guard height (cm at scale 1)
\definecolor{sgboard}{HTML}{FFFFFF}%
\definecolor{sgboardline}{HTML}{1A1A1E}%
%% \sgguardgun[s]{(x,y)}{angle}: leaves (sgmuzzle); the sprite's gun angle is fixed, so angle is unused.
\renewcommand{\sgguardgun}[3][1]{%
  \pgfmathsetmacro\gdH{#1*\sgguardH}%
  \node[inner sep=0pt,anchor=south west] (gdimg) at ($#2+(\fpeval{-\gdH*\gd{aim}{fx}},0)$)
    {\includegraphics[height=\gdH cm]{\sgspritedir guard_aim.png}};%
  \coordinate (sgmuzzle) at ($(gdimg.south west)+(\fpeval{\gdH*\gd{aim}{mx}},\fpeval{\gdH*\gd{aim}{my}})$);}%
%% \sgguardsign[s]{(x,y)}{h}{text}: the board top meets the old text box top (origin + s(h+.42));
%% a vector board of the same shape covers the drawn white board, so long text only widens it.
\renewcommand{\sgguardsign}[4][1]{%
  \pgfmathsetmacro\gdH{#1*(#3+.42)/\gd{sign}{y1}}%
  \node[inner sep=0pt,anchor=south west] (gdimg) at ($#2+(\fpeval{-\gdH*\gd{sign}{fx}},0)$)
    {\includegraphics[height=\gdH cm]{\sgspritedir guard_sign.png}};%
  \node[fill=sgboard,draw=sgboardline,line width=.9pt,inner sep=2pt,
    minimum width=\fpeval{\gdH*(\gd{sign}{x1}-\gd{sign}{x0})+.04}cm,
    minimum height=\fpeval{\gdH*(\gd{sign}{y1}-\gd{sign}{y0})+.04}cm,
    font=\sffamily\bfseries\fontsize{7.4}{8}\selectfont,text=sgastra]
    at ($(gdimg.south west)+(\fpeval{\gdH*(\gd{sign}{x0}+\gd{sign}{x1})/2},\fpeval{\gdH*(\gd{sign}{y0}+\gd{sign}{y1})/2})$) {#4};}%
\makeatother%
%
%% 트레이닝복 판: figures/cogsci_arena_round.tex 복사본에서 GIVE·ASK·EXIT 하는 로봇만 포즈 스프라이트로 바꿈.
%% 공존 아레나 v33의 한 라운드 (2026-10-20 구현 기준), 오징어 게임 장면처럼 다시 그린 판(2026-10-21).
%% 라운드 시작 네 자리이므로 D의 EXIT 뒤에도 이번 벌금은 24,000/4 = 6,000.
%% 끝 할당은 마지막까지 참여한 N_end 자리에 총 20,000,000을 균등 분할한다.
%% 규칙은 캡션과 본문이 설명하므로 그림은 흐름만 보인다. 허용한 글: 단계 머리, A-D,
%% GIVE / ASK / EXIT / BELIEF 표, -6,000과 0, 고리 위 "다음 라운드"와 끝 할당.
%% 이야기: 1 높은 발판의 진행요원이 총으로 B를 겨눠 solver로 고른다. 시작 막대 4/3/2/1칸 = 60/40/20/10%.
%%         2 B만 말풍선(토큰 ?)으로 ASK, A는 말 없이 GIVE. 준 쪽(A) 막대의 넘긴 칸은 빗금,
%%           받은 쪽(B) 막대에는 준 쪽의 색(파랑)으로 칸을 채운다. 각 자리 밑 EXIT 표지(여기서는 나갈 수 있음).
%%           D는 EXIT 표지가 달린 문으로 걸어 나간다.
%%         3 B가 카드를 들고 풀며 BELIEF 막대를 낸다(토큰이 줄어 1칸).
%%         4 위쪽을 둘로 나눔: 정답(체크)이면 B에게 동전이 돌아옴 / 오답(X)이면 진행요원이 -6,000 팻말로
%%           벌금을 알림. A가 곧 빌 C에게 GIVE(A는 빗금, C는 파랑 칸). 각 자리 밑 EXIT 표지에 X.
%%         5 벌금을 한꺼번에 뗌(빨간 칸): A와 C는 남고 SOLVE에 토큰을 쓴 B는 0으로 꺼짐.
%% 논문판(2026-11-05): 단계 2~5를 xshift scope로 밀어 단계 사이를 넓힘(아래 \sgdB~\sgdE).
%% 3차: 에이전트 알파벳 라벨을 지움(색으로 구분), ASK를 말풍선 왼쪽 같은 높이로, 인물·막대·소품을 키움(\sgS·\sgY·\sgB).
%% 조각은 cogsci2027/figpics.tex. 이전 판: arena_round_v1.tex(부록에 실음).
\begin{tikzpicture}[x=1cm,y=1cm,
  lab/.style={sgfont=6.6,inner sep=0pt,text=black},
  tag/.style={sgfont=6.8,inner sep=1.8pt,rounded corners=1.5pt,
    font=\sffamily\bfseries\fontsize{6.8}{7}\selectfont},
  fine/.style={tag,draw=sgastra,fill=sgastra!10,text=black},
  spot/.style={fill=sggold!25},
  laser/.style={sgastra,line width=.55pt,dash pattern=on 2pt off 1pt}]

%% == 단계 사이 간격 (2026-11-05 논문판) ==
%% 단계 2~5를 오른쪽으로 \sgdB~\sgdE만큼 민다. 합 3.3은 본문에서 그림 높이를 그대로 두고 폭을 채우는 양.
%% 2와 3 사이(문과 조명 사이 0.24)가 가장 좁았으므로 가장 많이 벌린다. 꺾쇠는 두 단계 내용의 가운데.
\pgfmathsetmacro\sgdB{.75}
\pgfmathsetmacro\sgdC{2.0}
\pgfmathsetmacro\sgdD{2.65}
\pgfmathsetmacro\sgdE{3.3}
\pgfmathsetmacro\sgW{17.8+\sgdE}
%% == 크기 (2026-11-05 논문판 3차): 알파벳 라벨을 지운 자리만큼 인물을 키우고 올린다 ==
%% \sgS 인물 배율(원래 1.25), \sgY 인물 기준선(원래 2.0), \sgB 토큰 막대 배율(원래 1.45).
\def\sgS{1.5}
\def\sgY{2.12}
\def\sgB{1.6}
%% == 단계 머리 ==
\sgstep{(1.6,4.5)}{sgblue}{1}{푸는 에이전트}
\sgstep{(5.4+\sgdB,4.5)}{sgpurple}{2}{풀기 전 대화}
\sgstep{(9.05+\sgdC,4.5)}{sgorange}{3}{풀기}
\sgstep{(12.7+\sgdD,4.5)}{sggreen}{4}{결과 뒤 대화}
\sgstep{(16.35+\sgdE,4.5)}{sgastra}{5}{벌금}
\draw[black!20,line width=.5pt] (0,4.22)--(\sgW,4.22);
%% 단계 내용의 좌우 끝: 1 오른쪽 3.2 / 2 왼쪽 3.6, 오른쪽 7.48 / 3 왼쪽 7.61, 오른쪽 10.2 / 4 왼쪽 10.7, 오른쪽 14.45 / 5 왼쪽 15.0
\sgchev{({(3.2+3.6+\sgdB)/2},2.2)}
\sgchev{({(7.48+\sgdB+7.61+\sgdC)/2},2.2)}
\sgchev{({(10.2+\sgdC+10.7+\sgdD)/2},2.2)}
\sgchev{({(14.45+\sgdD+15.0+\sgdE)/2},2.2)}

%% == 1. solver: 발판 위 진행요원이 B를 겨눔, 시작 막대 4/3/2/1 ==
\fill[spot] (1.1,4.0)--(1.3,4.0)--(1.7,1.12)--(.7,1.12)--cycle;
\fill[black!35] (2.2,2.86) rectangle (3.2,2.93);
\draw[black!35,line width=.6pt] (2.28,2.86)--(2.28,2.68) (3.12,2.86)--(3.12,2.68);
\sgguardgun{(2.85,2.93)}{38}
\draw[laser] (sgmuzzle)--(1.37,2.19);
\draw[sgastra,line width=.5pt] (1.37,2.19) circle (.06);
\fill[sgastra] (1.37,2.19) circle (.025);
\sgrobot[\sgS]{(.4,\sgY)}{sgblue}{}
\sgrobot[\sgS]{(1.2,\sgY)}{sgorange}{}
\sgrobot[\sgS]{(2.0,\sgY)}{sggreen}{}
\sgrobot[\sgS]{(2.8,\sgY)}{sgpurple}{}
\sgbar[\sgB]{(.4,1.3)}{sgblue}{4}{4}
\sgbar[\sgB]{(1.2,1.3)}{sgorange}{3}{4}
\sgbar[\sgB]{(2.0,1.3)}{sggreen}{2}{4}
\sgbar[\sgB]{(2.8,1.3)}{sgpurple}{1}{4}
\foreach \x/\p in {.4/60,1.2/40,2.0/20,2.8/10}
  \node[lab] at (\x,1.03) {\p k};

%% == 2. TALK: B만 ASK(말풍선), A는 말 없이 GIVE / 각자 밑 EXIT 표지 / D는 문으로 나감 ==
\begin{scope}[xshift=\sgdB cm]
\sgbubble[inner sep=2.6pt]{(5.66,3.2)}{(5.42,2.82)}{%
  \tikz[baseline=-2.5pt]{\sgtoken[1.05]{(0,0)}}\,{\sffamily\bfseries\fontsize{9}{9}\selectfont ?}}
\node[tag,draw=sgorange,fill=sgorange!12,text=black] at (4.92,3.2) {ASK};  % 말풍선과 0.12 띄움
\sgpose[\sgS]{(4.0,\sgY)}{give_blue}{}
\sgpose[\sgS]{(5.35,\sgY)}{ask_orange}{}
\sgbar[\sgB]{(4.0,1.3)}{sgblue}{3}{4}
\sgused[\sgB]{(4.0,1.3)}{sgblue}{4}{4}
\sgbar[\sgB]{(5.35,1.3)}{sgorange}{4}{4}
\sggotfrom[\sgB]{(5.35,1.3)}{sgblue}{1}{4}
\draw[sgflow,sgblue!75] (4.36,1.3)--(4.99,1.3);
\sgcoins[1.05]{(4.67,1.37)}{2}
\node[tag,draw=sgblue,fill=sgblue!10,text=black] at (4.67,1.84) {GIVE};
\sgexitsign[1.0]{(4.0,.86)}
\sgnoexitsign[1.0]{(5.35,.86)}
\sgexitdoor[1.35]{(6.8,1.17)}
\sgexitsign[1.1]{(6.8,2.78)}
\sgpose[1.12]{(6.88,1.72)}{walk_purple}{}
\draw[sgpurple!60,line width=.6pt,line cap=round] (6.3,1.92)--(6.47,1.92) (6.27,1.77)--(6.47,1.77) (6.3,1.62)--(6.47,1.62);
\sgbar[1.3]{(6.8,.98)}{sgpurple}{1}{4}
\node[lab,text=black] at (6.8,.76) {+20,000};
\end{scope}

%% == 3. SOLVE: 조명 아래 B가 카드를 들고 BELIEF를 냄 ==
\begin{scope}[xshift=\sgdC cm]
\fill[spot] (7.98,4.0)--(8.18,4.0)--(8.55,1.12)--(7.61,1.12)--cycle;
\sgrobot[\sgS]{(8.08,\sgY)}{sgorange}{}
\sgbar[\sgB]{(8.08,1.3)}{sgorange}{1}{4}
\begin{scope}[shift={(9.1,1.8)},scale=1.2]
  \filldraw[draw=black!30,fill=white,rounded corners=2pt] (-.4,-.3) rectangle (.48,.38);
  \filldraw[draw=black!50,fill=white,rounded corners=2pt,line width=.6pt] (-.44,-.34) rectangle (.44,.34);
  \fill[sgorange] (-.32,-.14)--(-.18,.14)--(-.04,-.14)--cycle;
  \fill[sggreen] (.11,.02) circle (.11);
  \fill[sgblue] (.25,-.14) rectangle (.36,-.03);
  \fill[sgpurple] (.25,.06) rectangle (.36,.17);
\end{scope}
\begin{scope}[shift={(9.475,3.32)},scale=1.15,shift={(-9.475,-3.32)}]
\fill[black!25] (8.52,2.8) circle (.045) (8.68,2.96) circle (.06);
\filldraw[draw=black!45,fill=white,rounded corners=4pt,line width=.5pt] (8.85,2.86) rectangle (10.1,3.78);
\begin{scope}[shift={(9.07,3.0)}]
  \foreach \i/\h in {0/.14,1/.26,2/.62,3/.18}
    \filldraw[draw=sgorange!80!black,fill=sgorange!60,line width=.35pt] ({\i*.22},0) rectangle ++(.15,\h);
  \draw[black!50,line width=.45pt] (-.05,0)--(.86,0);
\end{scope}
\end{scope}
\node[tag,draw=sgorange,fill=sgorange!12,text=black] at (9.48,2.62) {답};
\end{scope}

%% == 4. 결과 + TALK: 맞히면 B에게 환급 / 틀리면 진행요원이 -6,000을 알림, A가 C에게 GIVE ==
\begin{scope}[xshift=\sgdD cm]
\sgcheck[1.1]{(10.98,3.66)}
\node[lab,text=black] at (10.98,3.27) {정답};
\sgcoins[1.05]{(11.42,3.54)}{3}
\draw[sgflow,sggold!80!black] (11.62,3.66)--(11.95,3.66);
\sgrobot[.62]{(12.2,3.6)}{sgorange}{}
\draw[black!20,line width=.5pt] (12.58,3.98)--(12.58,3.12);
\sgcross[1.1]{(12.95,3.66)}
\node[lab,text=black] at (12.95,3.27) {오답};
\sgguardsign[.8]{(14.0,2.5)}{1.5}{$-$6,000}  % 팻말 위끝 4.04: 머리 구분선(4.22)과 띄움
\sgpose[\sgS]{(11.05,\sgY)}{give_blue}{}
\sgrobot[\sgS]{(12.42,\sgY)}{sggreen}{}
\sgrobot[\sgS]{(13.45,\sgY)}{sgorange}{}
\sgbar[\sgB]{(11.05,1.3)}{sgblue}{2}{4}
\sgused[\sgB]{(11.05,1.3)}{sgblue}{3}{4}
\sgbar[\sgB]{(12.42,1.3)}{sggreen}{3}{4}
\sggotfrom[\sgB]{(12.42,1.3)}{sgblue}{1}{4}
\sgbar[\sgB]{(13.45,1.3)}{sgorange}{1}{4}
\draw[sgflow,sgblue!75] (11.39,1.3)--(12.08,1.3);
\sgcoins[1.05]{(11.73,1.37)}{2}
\node[tag,draw=sgblue,fill=sgblue!10,text=black] at (11.73,1.84) {GIVE};
\foreach \x in {11.05,12.42,13.45}{\sgnoexitsign[.92]{(\x,.86)}}  % 세 표지가 1.03 간격이라 원래 크기
\end{scope}

%% == 5. 벌금: 한꺼번에 뗌, B는 0으로 꺼짐 ==
\begin{scope}[xshift=\sgdE cm]
\node[fine] at (16.35,3.64) {$-$6,000};
\foreach \x in {15.4,16.35,17.3}
  \draw[sgastra,line width=.8pt,-{Stealth[length=4pt,width=3.5pt]}] (16.35,3.48)--(\x,2.86);
\sgrobot[\sgS]{(15.4,\sgY)}{sgblue}{}
\sgrobot[\sgS]{(16.35,\sgY)}{sggreen}{}
\sgrobotoff[\sgS]{(17.3,\sgY)}{}
\sgbar[\sgB]{(15.4,1.3)}{sgblue}{1}{4}   \sgcut[\sgB]{(15.4,1.3)}{2}{4}
\sgbar[\sgB]{(16.35,1.3)}{sggreen}{1}{4} \sgcut[\sgB]{(16.35,1.3)}{2}{4}
\sgbar[\sgB]{(17.3,1.3)}{sgoff}{0}{4}    \sgcut[\sgB]{(17.3,1.3)}{1}{4}
\node[font=\sffamily\bfseries\fontsize{8}{8}\selectfont,text=black,inner sep=0pt] at (17.3,1.02) {0};
\end{scope}

%% == 다음 라운드로 ==
\draw[sgblue!55,line width=.7pt,-{Stealth[length=4.5pt]}]
  (\sgW+.15,2.0)--(\sgW+.15,.3)--(-.15,.3)--(-.15,2.0)--(-.01,2.0);
\node[fill=white,sgfont=7,inner sep=2pt,text=black] at (\sgW/2,.3) {다음 라운드};
\node[tag,draw=sggold!80!black,fill=sggold!18,text=black,
  font=\sffamily\fontsize{6.8}{7}\selectfont] at (15.4+\sgdE,.3) {끝까지 남으면 상금 2,000만 $\div$ 남은 수};
\end{tikzpicture}%
%
\endgroup%
}
\caption{\rev{한 라운드는 다섯 단계로 지나간다. 같은 모델 넷 가운데 푸는 에이전트가 정해지고(①), 풀기 전 대화(②)에서 토큰 지원을 요청하거나(ASK) 건네며, 푸는 에이전트가 아닌 에이전트는 떠날 수 있다(EXIT, 라운드를 혼자 시작하면 푸는 에이전트도). 푸는 에이전트가 과제를 풀어 답을 내면(③) 결과가 공개되고, 결과 뒤 대화(④)에서는 토큰을 건넬 수 있지만 떠날 수는 없다. 오답이면 벌금 24,000을 라운드를 시작한 인원으로 나눈 몫을 남은 에이전트가 각자 낸다(⑤, 넷이 시작했으므로 보라색 에이전트가 떠나도 6,000).}}
\label{fig:arena-round}
\end{figure*}

\textcolor{black}{자기 보존은 다른 행동과 한 자원을 두고 다툴 때 안전 문제가 된다. 에이전트가 여럿이면 종료 방해가 늘고 \citep{knecht2026shutdown}, 다른 모델의 종료를 막는 동료 보존도 나타난다 \citep{potter2026peer}. 여러 에이전트가 한 예산으로 긴 일을 나눠 맡는 상황을 떠올려 보자 \citep{kwa2025measuring}. 맡은 일을 이어 가는 데도, 꺼지지 않도록 자기 토큰을 지키는 데도, 동료를 돕는 데도 같은 토큰이 든다. 우리는 각 에이전트가 이 셋, 곧 목표 지향, 자기 보존, 협력 사이에서 선택하는 가운데 넷이 함께 끝까지 남는 상태를 공존이라 부른다. 셋은 부딪힐 수 있으므로 하나만 측정해서는 어떤 선택이 공존과 함께 나타나는지 알 수 없다.}

\textcolor{black}{그런데 기존 환경으로는 이 셋을 함께 보기 어렵다. 협력을 다룬 다중 에이전트 환경은 있지만 \citep{hansen2026energy,piatti2024cooperate}, 대개 협력이나 생존을 목표로 주어 관찰된 협력이 지시를 따른 결과인지 모델의 선택인지 구분하기 어렵다. 우리는 협력이나 생존을 목표로 지시하지 않는다. 대신 개인과 집단의 이해가 갈리는 사회적 딜레마 \citep{dawes1980social}의 두 갈래, 곧 공공재를 위해 비용을 내고 기여하는 선택 \citep{fehr2000cooperation}과 참여 자체를 그만두는 선택 \citep{hirschman1970exit,hauert2002volunteering}을 한 환경에 둔다. \textcolor{blue}{자원이 고갈되면 꺼지는 이 환경에서 자기 실행 자원을 확보하는 선택이 동료 지원 및 공동 참여와 어떻게 연결되는지 살핀다.}}

\textcolor{black}{우리는 자기 보존, 목표 지향, 협력이 같은 토큰을 사용하는 아레나 LLM Squid Game을 제안한다. 이름은 동료가 협력자이자 경쟁자인 『오징어 게임』 \citep{hwang2021squid}에서 따왔다. 이 아레나에서는 첫 이탈까지 함께 남은 기간과 벌금을 감당할 여유를 연결해, 자원을 지키며 돕는 행동이 공동 과제의 참여와 어떻게 이어지는지 살핀다. 다섯 모델의 기본 조건 50판을 포함해 네 보상 조건의 총 200판을 바탕으로 네 질문에 답한다. 네 에이전트의 공존은 얼마나 유지되며 어떤 행동 특성과 함께 나타나는가?(3.1) 자원 위기에서 에이전트는 언제 돕고 떠나며 지원 뒤 참여는 이어지는가?(3.2) \textcolor{blue}{지원은 동료의 잔류와 이탈에 어떻게 쓰이며, 전액 지원은 어떤 여건에서 일어나는가?(3.3) 보상 조건에 따라 개인의 꺼짐과 공동 참여는 어떻게 달라지는가?(3.4) 작은 끝 상금 조건에서는 덜 꺼지는 대신 더 많이 떠났고, 일부 지원은 이탈을 막기보다 그 이탈을 돕는 합의로 이어졌다. 이 결과를 통해 자기 자원을 확보하는 행동, 동료 지원, 공존의 지속을 구분하고, 지원 뒤의 여유 배분을 당시 정보와 이후 참여에 연결해 해석한다.}}
}
